\paragraph{Reflexión por clausura de testigos y construcción de la imagen.}
El paso de Colección--Reemplazo del lema~\ref{thm:ch-6b} admite la siguiente
demostración explícita. Se mantiene la jerarquía ya generada a partir del
código operatorio y del registro realizado; abreviamos
\(\mathfrak G=\mathfrak G_{\rm HMT}(h,m_\infty)\).
La argumentación utiliza Separación y Reemplazo en la metateoría ambiente,
sin presuponer estos esquemas en \(\mathfrak G\).

Sea \(\Phi\) una colección finita de fórmulas. Se traducen primero sus
cuantificadores universales mediante negación y cuantificación existencial,
y se cierra la colección resultante bajo subfórmulas. Para cualquier nivel
inicial existe un ordinal límite posterior \(\theta\) tal que
\[
 \mathfrak G_\theta\models\varphi(\bar a)
 \quad\Longleftrightarrow\quad
 \mathfrak G\models\varphi(\bar a)
 \qquad
 (\varphi\in\Phi,\ \bar a\in\mathfrak G_\theta^{<\omega}).
\]
La notación de satisfacción de la clase se interpreta, para cada fórmula
fijada, relativizando sus cuantificadores a la clase definible
\(\mathfrak G\); no se introduce un predicado uniforme de verdad.

\begin{proof}
Fijemos \(\alpha\). Para cada subfórmula existencial
\(\exists y\,\psi(y,\bar x)\) de \(\Phi\), las tuplas
\(\bar a\in\mathfrak G_\alpha^{<\omega}\) de la aridad correspondiente
que poseen un testigo en \(\mathfrak G\) forman un conjunto por Separación
metateórica. A cada una se asigna el menor ordinal \(\beta\) para el que
algún \(y\in\mathfrak G_\beta\) satisface
\(\mathfrak G\models\psi(y,\bar a)\). El ordinal existe porque
\(\mathfrak G\) es la unión de sus niveles. Reemplazo en la metateoría
acota esas primeras etapas de testigos. La finitud de \(\Phi\) permite
elegir definiblemente una sola cota \(f(\alpha)>\alpha\), válida para
todas sus subfórmulas existenciales. Por tanto, cada tupla considerada tiene
\emph{al menos un testigo} en \(\mathfrak G_{f(\alpha)}\).

Tomamos \(\alpha_0\) suficientemente grande para contener los parámetros
iniciales y ponemos
\[
 \alpha_{n+1}=f(\alpha_n),\qquad
 \theta=\sup_{n<\omega}\alpha_n.
\]
Por continuidad de la jerarquía, toda tupla finita de
\(\mathfrak G_\theta\) pertenece a algún \(\mathfrak G_{\alpha_n}\).
Si una fórmula existencial de \(\Phi\) con esa tupla es verdadera en
\(\mathfrak G\), posee un testigo en
\(\mathfrak G_{\alpha_{n+1}}\subseteq\mathfrak G_\theta\).
La inducción sobre las fórmulas prueba ahora la equivalencia: las atómicas
se conservan en la estructura inducida; los conectivos preservan la
equivalencia; el paso existencial utiliza ese testigo en una dirección y el
mismo testigo, visto en \(\mathfrak G\), en la otra. La traducción previa
de los universales incluye también las fórmulas existenciales que expresan
sus posibles contraejemplos. Queda probada la reflexión finita requerida.
\end{proof}

Sea ahora \(F(x,y,\bar p)\) funcional sobre \(a\) en \(\mathfrak G\).
Aplicamos la construcción a las fórmulas que expresan \(F\), la existencia
de sus testigos y su imagen, con \(a,\bar p\in\mathfrak G_\theta\).
La transitividad da \(a\subseteq\mathfrak G_\theta\). Cada
\(x\in a\) posee su valor dentro de ese nivel, y la reflexión identifica
la imagen exacta con
\[
 b=\{y\in\mathfrak G_\theta:
       \mathfrak G_\theta\models
       \exists x\in a\ F(x,y,\bar p)\}.
\]
Este conjunto es definible sobre \(\mathfrak G_\theta\), con parámetros
del mismo nivel. Por la operación sucesora de la jerarquía,
\(b\in\operatorname{Def}(\mathfrak G_\theta)
\subseteq\mathfrak G_{\theta+1}\). Así se obtiene Reemplazo en
\(\mathfrak G\). Si se conserva la existencia de testigos y se omite la
unicidad, la misma construcción proporciona un conjunto de testigos para
Colección.

\paragraph{Aridad de los nombres y selección interna.}
En el buen orden del lema~\ref{thm:ch-6b}, la enumeración de la base
asigna a cada elemento su menor código natural. En cada sucesor se ordenan
primero los códigos de fórmula; fijado uno, queda fijada la aridad de su
tupla de parámetros. El orden lexicográfico se aplica entonces a un
\emph{producto finito de aridad fija} de los buenos órdenes anteriores.
El menor nombre representa cada valor nuevo y los niveles límite conservan
el orden por primera aparición. De este modo queda especificado el orden
de las tuplas, sin mezclar aridades distintas en un único orden
lexicográfico. La selección del menor elemento de cada conjunto no vacío,
junto con Reemplazo ya obtenido, forma internamente las funciones de
elección utilizadas después. Esta construcción de nombres precede a las
inyecciones cardinales y no recibe ninguna biyección del continuo como
dato.
